\documentclass[11pt]{article}
\usepackage[margin=1in]{geometry}
\usepackage{xcolor}
\usepackage{booktabs}
\usepackage{tabularx}
\usepackage{array}
\usepackage[hidelinks]{hyperref}
\usepackage{listings}
\usepackage{fancyhdr}
\usepackage{enumitem}
\usepackage{amsmath,amssymb}
\usepackage{textcomp}

\definecolor{denim}{HTML}{2C6FBB}
\definecolor{denimdark}{HTML}{1B3A5C}
\definecolor{denimmid}{HTML}{3A6B8C}
\definecolor{denimrow}{HTML}{E6EEF5}
\definecolor{plum}{HTML}{7B2D8E}

\hypersetup{colorlinks=true, linkcolor=denim, urlcolor=denim, citecolor=denim}

\pagestyle{fancy}
\fancyhf{}
\fancyhead[L]{\footnotesize\color{denimmid}CS7641 SL Reproducibility --- Timothy Leung}
\fancyhead[R]{\footnotesize\color{denimmid}gtID 904150400 \,\textbar\, \texttt{tleung37@gatech.edu}}
\fancyfoot[C]{\thepage}
\renewcommand{\headrulewidth}{0.4pt}

\lstset{
  basicstyle=\ttfamily\footnotesize\color{denimdark},
  keywordstyle=\color{denim}\bfseries,
  commentstyle=\color{denimmid}\itshape,
  stringstyle=\color{plum},
  showstringspaces=false, breaklines=true, frame=single,
  rulecolor=\color{denimrow}, framesep=3pt, basewidth=0.5em,
}

\title{\color{denimdark}\textbf{Reproducibility Sheet --- CS7641 Supervised Learning Report}\\
       \large Covertype, five-learner comparison, Summer 2026}
\author{Timothy Leung \\ \texttt{tleung37@gatech.edu} \quad gtID 904150400 \\
        Georgia Institute of Technology, OMSCS}
\date{}

\begin{document}
\maketitle
\thispagestyle{fancy}

\section*{1.\ Overleaf project (read-only)}
\textbf{Overleaf project:} \href{https://www.overleaf.com/project/6a11cdb664b3e344a5bd3a25}{\texttt{overleaf.com/project/6a11cdb664b3e344a5bd3a25}}

The project is hosted on Georgia Tech Enterprise Overleaf
(\href{https://overleaf.gatech.edu}{overleaf.gatech.edu}); the read-only link
above grants the grader full revision-history access without write
permission. Revision timeline is intact from first commit through final
submission build.

\section*{2.\ GitHub commit hash (single SHA)}
\textbf{Final commit SHA:} \texttt{e287cb055afd8f20edb1646ad53574f57676adb8}\\
\textbf{Repository:} \href{https://github.gatech.edu/tleung37/CS7641_ml_report1}{\texttt{github.gatech.edu/tleung37/CS7641\_ml\_report1}} (GT Enterprise, private)\\
\textbf{File-link pattern:} \texttt{<repo>/blob/main/<filename>} --- directory layout is stable across revisions (only SHAs and minor edits change)\\
\textbf{Branch:} \texttt{main}\\
\textbf{Tag:} \texttt{v1.0-submission}

The repository contains every script that produces every figure, every
table, and every number in the report. Branch and commit history are intact
and available for review.

\section*{3.\ Exact run instructions (standard Linux machine)}

\subsection*{3.1\ Environment setup}
Tested on Ubuntu 22.04 LTS, Python 3.11, R 4.4.1. Approximate end-to-end
runtime: 25--30 minutes on a 4-core CPU laptop (no GPU required).

\begin{lstlisting}[language=bash]
# Clone the repository
git clone https://github.gatech.edu/tleung37/CS7641_ml_report1.git
cd CS7641_ml_report1
git checkout fb10855dc58c147f4da7492d6a9d21dcf2e008e3   # or: git checkout main

# Python environment
python3.11 -m venv .venv
source .venv/bin/activate
pip install -r requirements.txt

# R environment (for the optional sanity-check; not required to grade the SL report)
Rscript -e 'install.packages(c("caret","rpart","kknn","kernlab","nnet",
            "rmarkdown","bookdown","yaml","dplyr","ggplot2","knitr"),
            repos="https://cloud.r-project.org")'
\end{lstlisting}

\subsection*{3.2\ Commands to reproduce every figure and table}
\begin{lstlisting}[language=bash]
# (a) Run the full pipeline end-to-end. Reads configs/config.yaml,
#     writes results/{logs,tables,figures}/.
python run_pipeline.py --config configs/config.yaml

# (b) Compile the SL report PDF (twice for cross-references).
cd report
pdflatex -interaction=nonstopmode SL_Report.tex
pdflatex -interaction=nonstopmode SL_Report.tex
cd ..

# (c) Optional: rerun the R sanity check (independent reimplementation).
Rscript -e 'rmarkdown::render("rsanity/sanity_check.Rmd",
            output_format="bookdown::pdf_document2")'

# (d) Compile this reproducibility sheet.
cd repro
pdflatex -interaction=nonstopmode REPRO_SL.tex
\end{lstlisting}

\subsection*{3.3\ Data paths}
\texttt{sklearn.datasets.fetch\_covtype()} is called by
\texttt{src/data\_loader.py} on first run; the dataset (\textasciitilde\,75\,MB)
caches to \texttt{data/covtype\_cache.npz}. No manual download is required;
the script will fetch from UCI if no cache is present.

\subsection*{3.4\ Random seeds}
A single seed --- \texttt{seed = 7641} (the course number) --- is set globally
in \texttt{configs/config.yaml} and propagated to:
\texttt{numpy.random.seed}, \texttt{random.seed},  

\texttt{torch.manual\_seed}, \texttt{sklearn.utils.check\_random\_state},
the \texttt{random\_state} argument of every
\texttt{StratifiedShuffleSplit}/\texttt{StratifiedKFold}, and the
\texttt{set.seed()} call in 

\texttt{rsanity/sanity\_check.Rmd}. Re-running on
a clean checkout reproduces every reported number to within floating-point
precision on the same Python and \texttt{scikit-learn} versions; minor
\textasciitilde\,$10^{-3}$ drift may appear across BLAS implementations
(MKL vs.\ OpenBLAS) on the kernel SVM and MLP runs.

\subsection*{3.5\ Outputs landing locations}
\begin{itemize}[leftmargin=*,nosep]
  \item \texttt{results/logs/} --- run logs, leakage-probe outputs, timing,
        and the locked \texttt{HYPOTHESIS.md} hash.
  \item \texttt{results/tables/table\_grid\_*.csv} --- per-learner CV-grid
        results with fold mean and std (one CSV per learner).
  \item \texttt{results/tables/table\_per\_class\_*.csv} --- per-class
        precision/recall/F1 on the held-out 4{,}000-row test set.
  \item \texttt{results/tables/table\_final\_scores.csv} --- the headline
        leaderboard.
  \item \texttt{results/figures/*.png} --- class balance, learning curves,
        complexity curves, confusion matrices, runtime, NN epoch loss.
  \item \texttt{report/SL\_Report.pdf} --- the compiled report.
  \item \texttt{rsanity/sanity\_check.pdf} --- the optional R cross-check
        (independent implementation, not required for grading).
\end{itemize}

\section*{4.\ Full-data EDA summary}
The full UCI Covertype dataset (\texttt{sklearn.datasets.fetch\_covtype})
contains 581{,}012 rows, 54 features (10 continuous cartographic +
4 binary wilderness-area indicators + 40 binary soil-type indicators), and
seven cover-type classes. Class distribution and continuous-feature
summaries follow.

\begin{table}[h]\centering
\caption{Full Covertype class distribution. Class-frequency ratio across
the poles is $283{,}301 / 2{,}747 \approx 103{:}1$.}
\begin{tabularx}{0.85\linewidth}{Xrr}
\toprule
\textcolor{denimdark}{\textbf{Class}} & \textbf{Count} & \textbf{Share (\%)} \\
\midrule
1 Spruce/Fir          & 211{,}840 & 36.46 \\
2 Lodgepole Pine      & 283{,}301 & 48.76 \\
3 Ponderosa Pine      &  35{,}754 &  6.15 \\
4 Cottonwood/Willow   &   2{,}747 &  0.47 \\
5 Aspen               &   9{,}493 &  1.63 \\
6 Douglas-fir         &  17{,}367 &  2.99 \\
7 Krummholz           &  20{,}510 &  3.53 \\
\midrule
Total                 & 581{,}012 & 100.00 \\
\bottomrule
\end{tabularx}
\end{table}

\begin{table}[h]\centering
\caption{Full-data continuous-feature summary. Elevation is in metres,
Aspect/Slope in degrees, Hillshade columns on a 0--255 byte range, the three
horizontal-distance columns in metres.}
\begin{tabularx}{0.95\linewidth}{Xrrrr}
\toprule
\textcolor{denimdark}{\textbf{Feature}} & \textbf{min} & \textbf{max} & \textbf{mean} & \textbf{std} \\
\midrule
Elevation                   & 1859 & 3858 & 2959 & 280 \\
Aspect                      &    0 &  360 &  155 & 112 \\
Slope                       &    0 &   66 &   14 &   7 \\
Hillshade\_9am              &    0 &  254 &  212 &  27 \\
Hillshade\_Noon             &    0 &  254 &  223 &  20 \\
Hillshade\_3pm              &    0 &  254 &  142 &  38 \\
Horiz.\ Dist.\ Hydrology    &    0 & 1397 &  269 & 212 \\
Horiz.\ Dist.\ Roadways     &    0 & 7117 & 2350 & 1559 \\
Horiz.\ Dist.\ Fire Points  &    0 & 7173 & 1980 & 1325 \\
\bottomrule
\end{tabularx}
\end{table}

The 44 binary indicators (4 wilderness areas, 40 soil types) are
zero-or-one columns; each row has exactly one wilderness flag set and
exactly one soil-type flag set, so the effective categorical cardinality
is $4 \times 40 = 160$ wilderness-by-soil cells. The continuous columns
are not Gaussian (Elevation is bimodal at the two major timberline bands;
Hillshade\_3pm is right-skewed by deep-shadow zero values), which informs
the choice to use $z$-score standardisation rather than Box--Cox.

\section*{5.\ Sampled-data EDA summary}
The 20{,}000-row stratified subsample (used in every experiment) preserves
the full-data class proportions to within $\pm 0.01$ percentage points by
construction. Continuous-feature moments match the full data within $\pm
1\%$ on each of mean, std, min, and max --- the expected behaviour of
stratification on labels combined with representative random sampling on
features.

\begin{table}[h]\centering
\caption{Sampled-data ($n{=}20{,}000$) class distribution.}
\begin{tabularx}{0.85\linewidth}{Xrr}
\toprule
\textcolor{denimdark}{\textbf{Class}} & \textbf{Count} & \textbf{Share (\%)} \\
\midrule
1 Spruce/Fir          & 7{,}292 & 36.46 \\
2 Lodgepole Pine      & 9{,}751 & 48.76 \\
3 Ponderosa Pine      & 1{,}231 &  6.16 \\
4 Cottonwood/Willow   &      95 &  0.47 \\
5 Aspen               &     327 &  1.64 \\
6 Douglas-fir         &     598 &  2.99 \\
7 Krummholz           &     706 &  3.53 \\
\midrule
Total                 & 20{,}000 & 100.00 \\
\bottomrule
\end{tabularx}
\end{table}

\begin{table}[h]\centering
\caption{Sampled-data continuous-feature summary ($n{=}20{,}000$).}
\begin{tabularx}{0.95\linewidth}{Xrrrr}
\toprule
\textcolor{denimdark}{\textbf{Feature}} & \textbf{min} & \textbf{max} & \textbf{mean} & \textbf{std} \\
\midrule
Elevation                   & 1860 & 3858 & 2960 & 281 \\
Aspect                      &    0 &  360 &  155 & 112 \\
Slope                       &    0 &   65 &   14 &   7 \\
Hillshade\_9am              &   68 &  254 &  212 &  27 \\
Hillshade\_Noon             &    0 &  254 &  223 &  20 \\
Hillshade\_3pm              &    0 &  252 &  143 &  38 \\
Horiz.\ Dist.\ Hydrology    &    0 & 1383 &  269 & 211 \\
Horiz.\ Dist.\ Roadways     &    0 & 7023 & 2350 & 1559 \\
Horiz.\ Dist.\ Fire Points  &    0 & 6990 & 1980 & 1325 \\
\bottomrule
\end{tabularx}
\end{table}

A subsequent stratified 80/20 split of the 20{,}000-row subsample yields a
16{,}000-row training partition and a 4{,}000-row test partition that retain
the same class proportions; the held-out test set contains
$\{1459, 1950, 246, 19, 65, 120, 141\}$ examples per class respectively,
which is the per-class support reported alongside precision/recall/F1 in
the main report's per-class table.

\section*{6.\ Exact sampling procedure}

The 20{,}000-row working subsample is drawn from the 581{,}012-row UCI
Covertype dataset via a two-stage stratified protocol, both stages seeded
with \texttt{seed = 7641}:

\begin{enumerate}[leftmargin=*]
  \item \textbf{Stage 1 (sample).} A single-fold
        \texttt{StratifiedShuffleSplit(n\_splits=1, train\_size=20000,
        random\_state=7641)} draws \emph{exactly} 20{,}000 row-indices from
        the 581{,}012 full rows, stratified on the seven-class label
        column. Counts per class are preserved to within one row of the
        nearest integer realisation of full-data proportions; the resulting
        per-class counts are as in Table 4 above.
  \item \textbf{Stage 2 (held-out test split).} A second
        \texttt{StratifiedShuffleSplit(n\_splits=1, 
        
        test\_size=0.20,
        random\_state=7641)} on the 20{,}000-row sample produces a 16{,}000-row training partition and a 4{,}000-row test partition,
        again stratified on the label.
\end{enumerate}

\noindent\textbf{Standardisation contract.} The \texttt{StandardScaler} for
the ten continuous columns is fit \emph{only} on the 16{,}000-row training
partition and applied to both train and test. The 44 binary indicator
columns pass through unchanged. The whole pipeline is encoded in
\texttt{src/data\_loader.py} so no learner can accidentally see the test
partition during tuning, and no scaler statistic ever sees a test row.

\noindent\textbf{Tuning protocol.} Every learner's hyperparameter grid is
searched under stratified $k$-fold CV on the 16{,}000-row training
partition only; the test partition is sealed until the held-out evaluation
in \S\,V of the main report. Each of the five folds preserves the same
class proportions as the 20k sample. The SVM uses $k{=}3$ folds on a further
stratified 8{,}000-row subsample of the 16{,}000-row training partition
for tractability (RBF kernel evaluations scale super-linearly in $n$); the
final SVM model is refit on the full 16{,}000-row training partition with
the best hyperparameters.

\noindent\textbf{Leakage controls.} Three probes run before any final
number is believed: a stratified 

\texttt{DummyClassifier} (chance-floor
verification), a shuffled-label retrain of a depth-8

\texttt{DecisionTreeClassifier} (shortcut-information verification), and a
direct \texttt{numpy.intersect1d} on train/test index arrays at split time
(disjointness verification). Observed values:
$\text{Dummy}_{\text{acc}}{=}0.372$,
$\text{Dummy}_{\text{macro-F1}}{=}0.150$,
$\text{Shuffled-DT}_{\text{acc}}{=}0.484$,
$\text{Shuffled-DT}_{\text{macro-F1}}{=}0.106$,
$|\text{train} \cap \text{test}|{=}0$.
All three agree with the no-leakage null.

\section*{7.\ Hardware and software versions used to generate the submission}
\begin{itemize}[leftmargin=*,nosep]
  \item OS: Ubuntu 22.04 LTS (kernel 6.5)
  \item CPU: 4-core x86\_64; no GPU used (PyTorch on CPU device)
  \item Python: 3.11.6
  \item Key packages: \texttt{numpy 1.26.x}, \texttt{scikit-learn 1.4.x},
        \texttt{torch 2.2.x (CPU)}, \texttt{pandas 2.2.x},
        \texttt{matplotlib 3.8.x}, \texttt{pyyaml 6.x}
  \item R: 4.4.1; \texttt{caret 6.0}, \texttt{rpart 4.1}, \texttt{kknn 1.3.1},
        \texttt{kernlab 0.9}, \texttt{nnet 7.3}, \texttt{rmarkdown 2.27},
        \texttt{bookdown 0.40}
  \item LaTeX: TeX Live 2024 / TinyTeX with the IEEEtran conference
        document class (\texttt{IEEEtran.cls v1.8b})
\end{itemize}

\end{document}
